\begin{postsub}[Post-submission section]
This thesis has so far presented two routes to density regression in two different
languages: Section~\ref{sec:density regression} works with the log quantile density (LQD)
transformation in the Hilbert space $\mathcal{L}^2([0,1])$, whereas
Section~\ref{sec:fréchet regression} works intrinsically in the Wasserstein space
$(\mathcal{G}, d_W)$. In this section we place both methods in a single template, which
makes their similarities precise and gives a geometric explanation for the gap between
them observed in the simulation study of Section~\ref{sec:application}.

\subsection{A common template}
\label{sec:common_template}
Both methods share the structure
\begin{equation}
    \label{eq:transformation_template}
    \mathcal{G} \xrightarrow{\ \Phi\ } \mathcal{H}
    \xrightarrow{\ \text{linear regression}\ } \mathcal{H}
    \xrightarrow{\ \Phi^{-1}\ } \mathcal{G},
\end{equation}
that is, a \emph{representation map} (or chart) $\Phi$ sends a density to an element of a
Hilbert space $\mathcal{H}$, a linear regression is carried out there, and the fit is
mapped back to $\mathcal{G}$ for interpretation. The two methods differ only in the choice
of $\Phi$.

\paragraph{LQD method.} Here $\Phi = \psi$ is the LQD transformation
\eqref{eq:lqd_definition} and $\mathcal{H} = \mathcal{L}^2([0,1])$. The image is
unconstrained, so the regression reduces to coefficientwise ordinary least squares. The
price is twofold: $\psi$ is a nonlinear map, and it is \emph{blind to location}, since a
density $f$ and any translate $f(\cdot - c)$ share the same quantile density
(cf.\ Section~\ref{sec:transformation_interpretation}). The location must therefore be
carried as a separate scalar $c = Q(u^*)$ and reconstructed on inversion.

\paragraph{Fréchet method.} Here $\Phi = Q$ is the quantile representation and again
$\mathcal{H} = \mathcal{L}^2([0,1])$, but the image is restricted to the closed convex
cone $\mathcal{Q}$ of non-decreasing functions. By Lemma~\ref{lemma:dqeqdw} the map
$\Phi : (\mathcal{G}, d_W) \to (\mathcal{Q}, \norm{\cdot}_2)$ is an \emph{isometry}
\parencite[cf.][Chapter~2]{PanaretosZemel2020}. As shown in
Definition~\ref{def:global_frechet}, global Fréchet regression is then ordinary linear
regression on the quantile functions $Q_i$ inside this cone, followed by the
$\mathcal{L}^2$-projection onto $\mathcal{Q}$ of Algorithm~\ref{alg:quadprog}. Because the
location is encoded in $Q$ itself, no separate anchor is needed.

Table~\ref{tab:method_comparison} summarises the correspondence.

\begin{table}[h!]
    \centering
    \begin{tabular}{lll}
        \toprule
                                     & LQD method              & Fréchet method \\
        \midrule
        Representation $\Phi$        & $\psi = \log q$         & $Q$ (quantile function) \\
        Image space $\mathcal{H}$    & $\mathcal{L}^2([0,1])$  & cone $\mathcal{Q}\subset\mathcal{L}^2([0,1])$ \\
        Constraints in $\mathcal{H}$ & none                    & monotonicity (convex cone) \\
        Location information         & lost; anchor $c=Q(u^*)$ & retained in $Q$ \\
        Regression step              & coefficientwise OLS     & weighted mean of $Q_i$ \\
        Return to $\mathcal{G}$      & $\psi^{-1}$ and anchor  & projection onto $\mathcal{Q}$ \\
        Relation to $d_W$            & nonlinear distortion    & isometry \\
        \bottomrule
    \end{tabular}
    \caption[Comparison of the two transformation maps]{The LQD and Fréchet methods as two
    instances of the template \eqref{eq:transformation_template}, differing only in the
    representation map $\Phi$.}
    \label{tab:method_comparison}
\end{table}

\subsection{Why the geometry favours the Fréchet method}
\label{sec:geometry_argument}
The two charts are not interchangeable, because the error criterion is itself defined in
one of them. The integrated squared error reported in Section~\ref{sec:application} is
measured in $d_W$, and by Lemma~\ref{lemma:dqeqdw} this equals the squared $\mathcal{L}^2$
distance between quantile functions. The Fréchet method therefore regresses in
\emph{exactly} the chart in which the loss is the flat $\mathcal{L}^2$ metric. Since
$\Phi = Q$ is an isometry, a linear model in $Q$-space is consistent with the Wasserstein
geometry: minimising in-chart squared error coincides with minimising $d_W$ risk. The only
departure from a plain linear fit is the projection onto the cone $\mathcal{Q}$; but
projection onto a closed convex set in a Hilbert space is non-expansive, and the true
conditional quantile function $m_\oplus(x)$ lies in $\mathcal{Q}$, so
\begin{equation}
    \label{eq:nonexpansive}
    \big\| \Pi_{\mathcal{Q}}\bigl(\tilde{Q}(x)\bigr) - m_\oplus(x) \big\|_2
    \le \big\| \tilde{Q}(x) - m_\oplus(x) \big\|_2,
\end{equation}
where $\Pi_{\mathcal{Q}}$ denotes the projection onto $\mathcal{Q}$. The monotonicity
correction can thus only reduce the $d_W$ error: it is, in this sense, ``free''.

For the LQD method the situation is different. The chart $\Phi = \psi$ is a nonlinear,
non-isometric distortion of $(\mathcal{G}, d_W)$. A linear model that is optimal in
$\psi$-space --- that is, in $\mathcal{L}^2$ error on the log quantile densities --- is in
general \emph{not} optimal for $d_W$ risk, because $\psi$ does not preserve the metric. In
addition, the location component discarded by $\psi$ must be reinstated from the separately
estimated anchors $c(p)$ (here by linear interpolation), and any error in this
reconstruction enters the $d_W$ risk directly and does not vanish as the fit of the log
quantile densities improves.

This gives a geometric, rather than merely computational, explanation for the persistent
gap in Table~\ref{tab:mean_ise_values}. Even when the densities are observed directly, the
LQD risk is dominated by representation and location error, whereas the Fréchet risk
reflects only estimation error in the chart aligned with the loss. The density-estimation
step adds a comparable amount of noise to both methods but does not change this ordering.
In short, the Fréchet method performs better here not by an accident of tuning, but
because it linearises the problem in the same geometry in which performance is measured.
\end{postsub}
